\documentclass[11pt]{article}
\usepackage[a4paper,margin=27mm]{geometry}
\usepackage[T1]{fontenc}
\usepackage{lmodern,amsmath,amssymb,amsthm,mathtools,booktabs,microtype,xcolor,hyperref}
\hypersetup{colorlinks=true,linkcolor=blue!50!black,citecolor=blue!50!black,urlcolor=blue!50!black}
\newtheorem{theorem}{Theorem}[section]
\newtheorem{lemma}[theorem]{Lemma}
\newtheorem{proposition}[theorem]{Proposition}
\newtheorem{corollary}[theorem]{Corollary}
\theoremstyle{remark}\newtheorem{remark}[theorem]{Remark}
\newcommand{\Ind}{\operatorname{Ind}}
\newcommand{\Irr}{\operatorname{Irr}}
\newcommand{\Reg}{\operatorname{Reg}}
\newcommand{\St}{\operatorname{St}}
\newcommand{\ord}{\operatorname{ord}}
\newcommand{\PSL}{\operatorname{PSL}}
\newcommand{\PGL}{\operatorname{PGL}}
\newcommand{\SL}{\operatorname{SL}}
\newcommand{\F}{\mathbb F}
\newcommand{\C}{\mathbb C}
\newcommand{\Q}{\mathbb Q}
\newcommand{\one}{\mathbf 1}
\newcommand{\RH}{\mathrm{RH}}
\title{Optimal families of zero detectors for $\PSL_2(q)$\\[4pt]
\large A positive Hecke factorization in odd characteristic}
\author{Research note prepared for the ProveIt project}
\date{27 November 2026}
\begin{document}
\maketitle
\begin{abstract}
Let $L/k$ be a Galois extension of number fields with group $\PSL_2(q)$,
where $q\equiv1\pmod4$ is an odd prime power. We show that two explicitly
specified intermediate fields, of degrees $(q^2-1)/2$ and $q(q-1)/2$,
detect every zero of $\zeta_L$, without Artin holomorphy or any Riemann
hypothesis. A positive factorization into finite-order Hecke $L$-functions
gives an explicit multiplicity inequality at every point other than $1$.
The larger field degree is the least possible maximum degree of an
arbitrary spectrally detecting family, except at $q=9$, where the exact
optimum is $36$ and a separate positive identity supplies zero detection.
We also give a three-field identity valid for every odd $q$, removing
Artin holomorphy from a broader family of reductions. These results address
the odd-characteristic question in the inspected ProveIt article; they
do not prove nonabelian GRH. All finite exceptional calculations are
included in an exact standard-library Python verifier.
\end{abstract}
\tableofcontents

\section{Source question and status}
The inspected ProveIt snapshot is
\texttt{ce35f320bc88c02093e479f5fe5003ff86b2795f}.
Its consolidated article, in \texttt{Article/sections/grp.tex}, asks in
Question \texttt{grp:q:odd} for the least possible maximum degree in a
spectrally detecting family for $\PSL_2(q)$ in odd characteristic, and asks
whether the Artin-holomorphy assumption can be removed from such
reductions. It already determines the least degree for a \emph{single}
spectrally detecting field and provides unconditional analytic transfers
in even characteristic. This note settles the maximum-degree question for
$q\equiv1\pmod4$, including its unique exception, and proves unconditional
analytic transfers. A supplementary three-field identity works for all
odd $q$.

The finite-group input is the ordinary character table of $\PSL_2(q)$,
obtained from the central-trivial rows of the $\SL_2(q)$ table in
\cite[Table 1]{BFL}, and Dickson's subgroup classification in
\cite[Theorem A.1]{GZ}. The analytical input is Artin formalism and the
classical fact that a linear Artin $L$-function is a finite-order Hecke
$L$-function. The argument uses none of the zero-location conclusions of
RH or Dirichlet GRH. Novelty is claimed relative to the inspected project
and is subject to an independent literature and proof review. No Lean
formalization or resolution of nonabelian GRH is claimed.

\section{The two fields and the main theorem}
Throughout Sections 2--6, let $q\equiv1\pmod4$ be an odd prime power,
$q\ge5$, and put
\[
 G=\PSL_2(\F_q),\qquad |G|=\frac{q(q^2-1)}2,\qquad
 a=\frac{q-1}2,\quad b=\frac{q+1}2.
\]
Let $B$ be the stabilizer of a point of $\mathbb P^1(\F_q)$, let $U$ be
its unipotent subgroup, and let $T_s=B/U$ and $T_n$ be respectively the
split and nonsplit projective tori. Their orders are
\[
 |U|=q,\quad |B|=q(q-1)/2,\quad |T_s|=a,\quad |T_n|=b.
\]
Write $N=N_G(T_n)$; it is dihedral of order $q+1=2b$. For a subgroup $H$
write $\pi_H=\Ind_H^G\one_H$. For a Galois extension $L/k$ with group
$G$, put $F_H=L^H$.

\begin{theorem}[Unconditional two-field zero detection]\label{thm:analytic}
For every such $L/k$ and every $s_0\in\C\setminus\{1\}$,
\begin{equation}\label{eq:orders}
 0\le\ord_{s_0}\zeta_{F_U},\ \ord_{s_0}\zeta_{F_N}
 \le\ord_{s_0}\zeta_L
 \le \frac{q+1}2\ord_{s_0}\zeta_{F_U}
       +(q-1)\ord_{s_0}\zeta_{F_N}.
\end{equation}
The middle notation means each of the two displayed subfield orders is at
most the order of $\zeta_L$. Consequently the zero set of $\zeta_L$ away
from $1$ is the union of the zero sets of these two subfield zeta functions,
and
\[
 \RH(\zeta_L)\quad\Longleftrightarrow\quad
 \RH(\zeta_{F_U})\ \text{and}\ \RH(\zeta_{F_N}).
\]
The field degrees are
\[
 [F_U:k]=\frac{q^2-1}2,\qquad [F_N:k]=\frac{q(q-1)}2.
\]
The quotient
\begin{equation}\label{eq:quotient}
 \frac{\zeta_{F_U}(s)^{(q+1)/2}\zeta_{F_N}(s)^{q-1}}{\zeta_L(s)}
\end{equation}
is holomorphic away from $1$, where it has a pole of exact order
$3(q-1)/2$.
\end{theorem}

A family $(F_{H_i})$ is \emph{spectrally detecting} if each nontrivial
irreducible character of $G$ occurs in at least one $\pi_{H_i}$.
This is the definition used by the source article. It is a finite-group
condition; it alone does not prevent zeros and poles of meromorphic Artin
factors from cancelling. Theorem~\ref{thm:analytic} supplies the extra
analytic assertion for our family.

\begin{theorem}[Optimal maximum degree]\label{thm:optimal}
Let $q\equiv1\pmod4$ be an odd prime power. Among all finite spectrally
detecting families of subfields of a $\PSL_2(q)$-extension, the smallest
possible value of the largest degree is
\begin{equation}\label{eq:opt}
 \mathcal D(q)=
 \begin{cases}
 36,&q=9,\\
 (q^2-1)/2,&q\ne9.
 \end{cases}
\end{equation}
For $q\ne9$, the pair $(F_U,F_N)$ attains the optimum and satisfies
Theorem~\ref{thm:analytic}. For $q=9$, a nonsplit torus normalizer $N$ of
order $10$ and a subfield subgroup $A\simeq\PSL_2(3)\simeq A_4$ of order
$12$ attain it. Adding the base field $k$, of degree one, gives an
unconditional zero-detecting family of the same optimal maximum degree.
\end{theorem}

The optimality in this theorem is among spectrally detecting families, not
a lower bound for every imaginable analytic criterion. In particular,
positivity of the Hecke factorization is an additional theorem rather than
an assumed consequence of spectral coverage.

\section{The character identity}
Let $z$ be the involution of $T_s$ and let $\delta$ be its unique quadratic
character. The principal series consists of
$P_\alpha=\Ind_B^G\alpha$, where $\alpha$ ranges over nontrivial
characters of $T_s$ other than $\delta$, with $\alpha$ and $\alpha^{-1}$
identified; these have degree $q+1$. The cuspidal series $D_\beta$ is
indexed by nontrivial characters of $T_n$, up to inversion; these have
degree $q-1$. The remaining characters are $\one$, the Steinberg character
$\St$ of degree $q$, and two exceptional characters $E_+,E_-$ of degree
$(q+1)/2$. Their sum is
\[
 E:=E_++E_-=\Ind_B^G\delta.
\]
Set
\[
 P_0=\sum_{\substack{\alpha\ne1,\delta\ ({\rm mod\ inversion})\\
                         \alpha(z)=1}}P_\alpha,
 \qquad
 I=\frac{1+\delta(z)}2=
 \begin{cases}1,&q\equiv1\pmod8,\\0,&q\equiv5\pmod8.\end{cases}
\]
The following table lists fixed-space dimensions, i.e. multiplicities in
the permutation characters. The subsequent analytic argument permits
negative orders for individual irreducible Artin functions.
\begin{center}
\begin{tabular}{@{}lrrrr@{}}\toprule
irreducible & degree & $\langle\pi_U,\chi\rangle$
 & $\langle\pi_N,\chi\rangle$ & $\langle\pi_B,\chi\rangle$\\\midrule
$\one$ & $1$ & $1$ & $1$ & $1$\\
$\St$ & $q$ & $1$ & $1$ & $1$\\
$P_\alpha$ & $q+1$ & $2$ & $1+\alpha(z)$ & $0$\\
$D_\beta$ & $q-1$ & $0$ & $1$ & $0$\\
$E_+$ & $(q+1)/2$ & $1$ & $I$ & $0$\\
$E_-$ & $(q+1)/2$ & $1$ & $I$ & $0$\\\bottomrule
\end{tabular}
\end{center}

\begin{lemma}[Fixed-space calculations]\label{lem:fixed}
The table is correct. In particular, $\pi_U+\pi_N$ contains every
irreducible character.
\end{lemma}
\begin{proof}
For reference, an ordinary principal character has value $1$ at a
nonidentity unipotent, value $\alpha(t)+\alpha(t)^{-1}$ at a split
semisimple element $t$, and value zero at a nonsplit semisimple element.
An ordinary cuspidal character has the values $-1$, zero, and
$-\beta(t)-\beta(t)^{-1}$ on these respective types. Steinberg has values
$0,1,-1$. Each exceptional character has value $\delta(t)$ at a split
semisimple element and zero at a nonsplit element; its two unipotent
values are $(1+\sqrt q)/2$ and $(1-\sqrt q)/2$, interchanged for the
other exceptional character. These are the central-trivial rows of
\cite[Table 1]{BFL}.

The group $U$ has $(q-1)/2$ nonidentity elements of each unipotent type.
Averaging gives the $U$ column. The two unipotent values of each
exceptional character sum to one, so its $U$-fixed dimension is
$((q+1)/2+(q-1)/2)/q=1$.

Since $b$ is odd, all $b$ involutions in $N\setminus T_n$ are split.
Each principal character is zero on $T_n\setminus\{1\}$ and has value
$2\alpha(z)$ on those involutions. Its average is therefore
$(q+1+2b\alpha(z))/(2b)=1+\alpha(z)$.
For each cuspidal parameter, $\sum_{t\ne1}\beta(t)=-1$, so its sum over
$T_n$ is $q-1+2=q+1$; the reflections contribute zero. This gives fixed
dimension one. The exceptional average is
$((q+1)/2+b\delta(z))/(2b)=I$.
Steinberg's average is $(q-(b-1)+b)/(2b)=1$.
Finally $\pi_B=\one+\St$ because the projective-line action is doubly
transitive. The table and the coverage assertion follow.
\end{proof}

\begin{proposition}[Positive character identity]\label{prop:identity}
In the character ring of $G$,
\begin{equation}\label{eq:identity}
 (q-1)\pi_N+\frac{q+1}2\pi_U-\Reg_G
 =(q-1)\left(\one+\frac12\pi_B+2P_0+I E\right).
\end{equation}
Every coefficient on the right is a nonnegative integer after the outer
factor is multiplied in, and every summand is induced from a linear
character of a subgroup.
\end{proposition}
\begin{proof}
Compare irreducible multiplicities using Lemma~\ref{lem:fixed}.
For a principal character the difference on the left is
$(q-1)(1+\alpha(z))$, which is $2(q-1)$ for the parameters in $P_0$ and
zero for the others. It is zero on every cuspidal character and
$(q-1)I$ on each exceptional character. The Steinberg coefficient is
$(q-1)/2$, and the trivial coefficient is $3(q-1)/2$.
These are exactly the coefficients on the right.
The principal series and $E$ are induced from nontrivial linear
characters of $B$. The character $\pi_B$ is induced from its trivial
character, and $\one$ is linear on $G$ itself.
\end{proof}

\section{From the identity to zeros}
For a point $s_0\ne1$ define the Heilbronn virtual character
\[
 \Theta=\sum_{\chi\in\Irr(G)}
       \ord_{s_0}L(s,\chi,L/k)\,\chi.
\]
Artin meromorphy gives meaning to the orders, including negative orders
when an individual Artin function has a pole. Artin induction and Hecke
holomorphy imply
\begin{equation}\label{eq:Heilbronn}
 \langle\Theta,\Ind_H^G\lambda\rangle
 =\ord_{s_0}L(s,\lambda,L/F_H)\ge0
\end{equation}
for every linear character $\lambda$ of every subgroup $H$. The trivial
case is a Dedekind zeta function, which also has no pole away from $1$.

Put $m=\Theta(1)=\ord_{s_0}\zeta_L$ and
$a_H=\langle\Theta,\pi_H\rangle=\ord_{s_0}\zeta_{F_H}$.
On each cyclic subgroup the restriction of $\Theta$ is an actual
character, by \eqref{eq:Heilbronn}; hence $|\Theta(g)|\le m$ for all
$g\in G$. Averaging over $H$ shows $0\le a_H\le m$ for every $H$.
Pairing \eqref{eq:identity} with $\Theta$ proves the last inequality in
\eqref{eq:orders}. This proves all zero-set and RH assertions of
Theorem~\ref{thm:analytic}.

More explicitly, Artin formalism gives the exact meromorphic identity
\begin{align}\label{eq:factorization}
 \frac{\zeta_{F_U}^{(q+1)/2}\zeta_{F_N}^{q-1}}{\zeta_L}
 &=\zeta_k^{q-1}\zeta_{F_B}^{(q-1)/2}
   \prod_{\substack{\alpha\ne1,\delta\ ({\rm mod\ inversion})\\
                    \alpha(z)=1}}
        L(s,\alpha,L/F_B)^{2(q-1)}
   L(s,\delta,L/F_B)^{I(q-1)}.
\end{align}
This includes the ramified Euler factors. The finite-order Hecke
characters in the product are nontrivial, so their $L$-functions are
entire and nonzero at $1$. The two zeta factors have simple poles there.
This proves the asserted holomorphy and the pole order $3(q-1)/2$.

\begin{remark}[What this does and does not prove]
The multiplicity bound is unconditional and holds over arbitrary number
fields. It prevents cancellation from concealing every off-line zero
from the specified fields. It does not establish RH for either smaller
field and consequently does not establish RH for $\zeta_L$. The source
Dirichlet-GRH theorem does not by itself provide RH for these nonabelian
intermediate fields.
\end{remark}

\section{Optimality for \texorpdfstring{$q\ne9$}{q different from 9}}
For a finite group $G$ define
\[
 h_{\rm fam}(G)=\max_{(H_i)}\min_i|H_i|,
\]
where the maximum ranges over spectrally detecting finite families.
Then the least possible maximum field degree is $|G|/h_{\rm fam}(G)$.
Lemma~\ref{lem:fixed} gives $h_{\rm fam}(G)\ge q$.
For $q\ge13$, and also for the obstruction calculation at $q=9$, choose
a faithful character $\alpha$ of the cyclic group $T_s$ and let
$P=\Ind_B^G\alpha$. It is irreducible of degree $q+1$ because
$\alpha\ne1,\delta$.

\begin{lemma}[The faithful principal-series obstruction]\label{lem:obstruction}
If $q\equiv1\pmod4$ and $q\ge13$, then $P^H=0$ for every subgroup
$H\le G$ with $|H|>q$.
\end{lemma}
\begin{proof}
We apply Dickson's classification \cite[Theorem A.1]{GZ}; $H=G$ is
immediate because $P$ is nontrivial irreducible.

\emph{Borel subgroups and their subgroups.}
Conjugate $H$ into $B=\F_q\rtimes(\F_q^\times)^2$, write
$A=H\cap U$, and let $C$ be its multiplier image.
If $C=1$ then $|H|\le q$. Otherwise translate the fixed point of a lift
of a generator of $C$ to zero, so that $H=A\rtimes C$.
For the subfield $\F_r$ generated by $C$, the additive group $A$ is an
$\F_r$-vector subspace. If $A\ne U$, then
$|A|\le q/r$ and $|C|\le r-1$, giving $|H|<q$.
Thus $|H|>q$ forces $A=U$ and $C\ne1$.
The sum of $P$ on $U$ is $(q+1)+(q-1)=2q$.
For each $c\in C\setminus\{1\}$, its $q$ affine lifts are all split
semisimple and contribute $q(\alpha(c)+\alpha(c)^{-1})$.
Faithfulness gives $\sum_{c\ne1}\alpha(c)=-1$ and the same for
$\alpha^{-1}$. The total sum on $H$ is therefore $2q-2q=0$.

\emph{Torus normalizers.}
A subgroup of a split torus normalizer has order at most $q-1$.
The only subgroup of the nonsplit normalizer of order greater than $q$
is the entire group $N$, of order $q+1$.
Since $a$ is even and $\alpha$ is faithful, $\alpha(z)=-1$.
Lemma~\ref{lem:fixed} gives $\dim P^N=0$.

\emph{Subfield subgroups.}
Suppose $H\simeq\PSL_2(t)$ with $q=t^j$, $j\ge2$, or
$H\simeq\PGL_2(t)$ with even $j\ge2$.
If $j\ge3$ in the first case or $j\ge4$ in the second, the subgroup
order is less than $q$. Thus only $q=t^2$ need be considered.
Here $t\ge5$. Every semisimple element of either subfield group is split
in $G$; both types of maximal torus have nontrivial faithful $\alpha$
restriction. There are $t(t+1)/2$ split tori and $t(t-1)/2$ nonsplit tori
in the subfield group, and each nonidentity semisimple element belongs
to exactly one such torus. In $\PSL_2(t)$ their orders are
$(t-1)/2,(t+1)/2$, and in $\PGL_2(t)$ their orders are $t-1,t+1$.
The sum of $P$ over the nonidentity elements of each torus is $-2$.
There are $t^2-1$ nonidentity unipotents, at which $P=1$.
Consequently the sum of $P$ over $H$ is
\[
 (t^2+1)+(t^2-1)-2\left(\frac{t(t+1)}2+
                                 \frac{t(t-1)}2\right)=0.
\]
This calculation is independent of any possible subgroup fusion into
the ambient group, because it uses complete cyclic tori.

\emph{The exceptional groups.}
The remaining groups are $A_4,S_4,A_5$, of orders at most $60$.
The condition $|H|>q\ge13$ leaves only the pairs in the following table.
Their existence conditions are the ones in Dickson's list; in each case
the character average is zero:
\begin{center}
\begin{tabular}{@{}rrrrrrr@{}}\toprule
$q$&$17$&$25$&$29$&$41$&$49$\\
$H$&$S_4$&$A_5$&$A_5$&$A_5$&$A_5$\\
$\sum_{h\in H}P(h)$&$18-18$&$26-30-20+24$&$30-30$
 &$42-30-12$&$50-30-20$\\\bottomrule
\end{tabular}
\end{center}
For completeness, $A_4,S_4,A_5$ have respectively
$(3,8,0,0)$, $(9,8,6,0)$, $(15,20,0,24)$ elements of orders
$(2,3,4,5)$. An element of order $d$ prime to the characteristic is
split precisely when $d\mid a$; the average of the faithful principal
character over all generators of a split cyclic subgroup of order $d$
is $2\mu(d)/\varphi(d)$. On a nonsplit subgroup it is zero, and on a
unipotent it is one. These rules give all entries exactly, with no
numerical approximation. They finish the proof.
\end{proof}

If every subgroup in a family had order greater than $q$, it would miss
the same character $P$. This proves $h_{\rm fam}(G)\le q$ for $q\ge13$.
Together with the upper construction, this proves
Theorem~\ref{thm:optimal} in those degrees.

\paragraph{The field size $q=5$.}
Here $G\simeq A_5$. Take either irreducible character $E$ of degree three.
The proper subgroups of order greater than five have types $A_4,D_{10},
D_6$. At an involution $E=-1$, and at an element of order three $E=0$.
The sum on the four nonidentity elements of a $C_5$ is two, because the
two five-cycle values are $(1\pm\sqrt5)/2$, each appearing twice.
Thus the sums over these three subgroups are respectively
$3-3=0$, $3+2-5=0$, and $3-3=0$.
The full group also has no invariant. Therefore no family using only
subgroups of order greater than five is spectrally detecting.
The pair $U,N$ has orders five and six, and establishes
$h_{\rm fam}(A_5)=5$, or $\mathcal D(5)=12$.

\section{The unique exception \texorpdfstring{$q=9$}{q = 9}}\label{sec:q9}
For $G=\PSL_2(9)\simeq A_6$, let $A\simeq\PSL_2(3)\simeq A_4$ be a
subfield subgroup and let $N$ again be a nonsplit torus normalizer, of
order ten. In an ordering of the seven irreducibles, their permutation
multiplicities are
\begin{center}
\begin{tabular}{@{}lrrrrrrr@{}}\toprule
 & $\one$&$\St$&$P$&$D_1$&$D_2$&$E_+$&$E_-$\\
degree&1&9&10&8&8&5&5\\\midrule
$\pi_N$&1&1&0&1&1&1&1\\
$\pi_A$&1&1&1&0&0&2&0\\
$\pi_B$&1&1&0&0&0&0&0\\
$\pi_{H_+}$&1&0&0&0&0&1&0\\
$\pi_{H_-}$&1&0&0&0&0&0&1\\\bottomrule
\end{tabular}
\end{center}
Here $H_+,H_-$ are subgroups isomorphic to $A_5$, of index six.
An index-six permutation character is $\one$ plus a degree-five
irreducible; the outer automorphism induced by $\PGL_2(9)$ swaps the
two unipotent classes and the two degree-five characters, and supplies
both choices. We label them so that $E_+$ occurs in $\pi_A$.

The $A$ row follows by averaging: its three involutions are split, its
eight elements of order three are unipotent in the same ambient class,
and the exceptional values on that class are $2,-1$.
Thus $\pi_N+\pi_A$ contains every irreducible, and the corresponding
field degrees are $36$ and $30$.

To prove optimality, take either degree-eight cuspidal character $D$.
Every subgroup of order greater than ten contained in a Borel contains
$U$, on which $D$ has no fixed vector, by the affine argument above.
No torus normalizer has order greater than ten. The remaining proper
subgroups of order greater than ten have types $A_4,S_4,A_5$.
For the first two, the sum of $D$ is $8-8=0$, since their elements of
order three are unipotent and the other nonidentity elements are split.
For $A_5$ the sum is $8-20+24/2=0$: elements of order five are nonsplit,
and the average cuspidal value on their generators is $1/2$.
The full group also fixes nothing. Hence $h_{\rm fam}(G)=10$ and
$\mathcal D(9)=36$.

The displayed multiplicities also prove the integral character identity
\begin{equation}\label{eq:q9}
 8\pi_N+10\pi_A+18\one-\Reg_G
 =9\pi_B+23\pi_{H_+}+3\pi_{H_-}.
\end{equation}
It gives the exact zeta identity
\[
 \frac{\zeta_{F_N}^{8}\zeta_{F_A}^{10}\zeta_k^{18}}{\zeta_L}
 =\zeta_{F_B}^{9}\zeta_{F_{H_+}}^{23}\zeta_{F_{H_-}}^{3}
\]
and, away from $1$, the inequality
\[
 \ord\zeta_L\le8\ord\zeta_{F_N}+10\ord\zeta_{F_A}+18\ord\zeta_k.
\]
As before each subfield zero is a zero of $\zeta_L$. The family
$(k,F_N,F_A)$ is therefore an unconditional zero detector with maximum
degree $36$. Over $\Q$, if RH for $\zeta$ is supplied as a separate
premise, the degree-one member adds no further RH assumption.
No zero-location premise is needed for the unconditional equivalence
with the three stated fields.

\section{A three-field identity for every odd \texorpdfstring{$q$}{q}}
In this section let $q\ge3$ be any odd prime power. Keep the same
definitions of $U,B,T_n$; the normalizer $N$ is not needed.

\begin{proposition}[All odd field sizes]\label{prop:allodd}
For $G=\PSL_2(q)$,
\begin{equation}\label{eq:allodd}
 \Reg_G=\frac{q-1}2\bigl(\pi_{T_n}+\pi_B-2\one\bigr)+\pi_U.
\end{equation}
Consequently, for every Galois $L/k$ with that group,
\[
 \zeta_L\zeta_k^{q-1}
 =\zeta_{F_{T_n}}^{(q-1)/2}\zeta_{F_B}^{(q-1)/2}\zeta_{F_U}.
\]
The three fields $F_{T_n},F_B,F_U$ detect all zeros of $\zeta_L$
away from $1$, and the RH for $\zeta_L$ is equivalent to their three
Riemann hypotheses without assuming Artin holomorphy. Their largest
degree is $q(q-1)$.
\end{proposition}
\begin{proof}
Every ordinary principal or cuspidal irreducible has two $T_n$-fixed
vectors, while Steinberg and each exceptional constituent have one.
The trivial character also has one. The $U$-fixed dimensions are one
for Steinberg, two for each ordinary principal character, zero for
each ordinary cuspidal character, and either one or zero for the
exceptional pair according as $q\equiv1$ or $3\pmod4$.
The exceptional degrees are correspondingly $(q+1)/2$ or $(q-1)/2$.
Comparing these multiplicities with the degrees gives
\eqref{eq:allodd}; it also holds for $q=3$, where the exceptional pair
is the pair of nontrivial linear characters of $A_4$.

Artin formalism gives the zeta identity. At $s_0\ne1$ all Dedekind
orders on its right, and $\ord\zeta_k$, are nonnegative, so
\[
 \ord\zeta_L\le\frac{q-1}2
      (\ord\zeta_{F_{T_n}}+\ord\zeta_{F_B})+\ord\zeta_{F_U}.
\]
Every subfield order is at most $\ord\zeta_L$ by the local positivity
argument in Section 4. This proves zero detection and the RH equivalence.
The degrees are $q(q-1)$, $q+1$, and $(q^2-1)/2$ respectively.
\end{proof}

This proposition does not claim optimality for $q\equiv3\pmod4$.
Its contribution is an unconditional transfer valid in all odd
characteristics; the sharper optimal theorem above addresses the other
congruence class.

\section{Example, verification, and precise remaining scope}
For $q=13$, the group has order $1092$. Theorem~\ref{thm:analytic}
uses fields of degrees $84$ and $78$, and gives
\[
 \ord_{s_0}\zeta_L\le7\ord_{s_0}\zeta_{F_U}
                       +12\ord_{s_0}\zeta_{F_N}\qquad(s_0\ne1).
\]
No spectrally detecting family can have all field degrees below $84$.
The source article's optimal single field for this value of $q$ has
degree $156$. If $\alpha$ denotes either primitive cubic character of
$B/U\simeq C_6$, the positive factorization is particularly short:
\[
 \frac{\zeta_{F_U}^{7}\zeta_{F_N}^{12}}{\zeta_L}
   =\zeta_k^{12}\zeta_{F_B}^{6}L(s,\alpha,L/F_B)^{24}.
\]
The inverse characters give the same induced Artin character, so choosing
either representative gives the same function.

The accompanying \texttt{verify\_psl2.py} uses only exact integers and
fractions. It checks the main identity in both class values and
irreducible multiplicities for the $94$ prime powers
$q\equiv1\pmod4$ with $5\le q\le1000$: $41{,}636$ class cases and
$21{,}006$ irreducible cases. It checks the all-odd identity at $183$
prime powers, all eight exceptional subgroup cases used in the
optimality proof (including the three $q=9$ cases), and the complete
$q=9$ identity. These finite checks confirm the algebra; the all-degree
theorems rest on the proofs above, not extrapolation from computations.

The paper resolves the inspected odd-characteristic question for the
maximum degree when $q\equiv1\pmod4$, and removes the Artin-holomorphy
assumption by explicit positive identities. The exact optimal
maximum-degree classification when $q\equiv3\pmod4$ is not proved here.
Neither RH for the smaller nonabelian fields nor a prominent general
conjecture such as Artin holomorphy or nonabelian GRH is proved.

\begin{thebibliography}{9}
\bibitem{BFL}
B. B\"ohmler, N. Farrell, and C. Lassueur,
\emph{Trivial source character tables of $\mathrm{SL}_2(q)$},
Journal of Algebra \textbf{598} (2022), 308--350,
doi:10.1016/j.jalgebra.2022.01.024.
\url{https://arxiv.org/abs/2107.06241}. See Table 1.
\bibitem{GZ}
R. M. Guralnick and M. E. Zieve,
\emph{Polynomials with $\mathrm{PSL}(2)$ monodromy},
Annals of Mathematics \textbf{172} (2010), 1315--1359,
doi:10.4007/annals.2010.172.1315.
\url{https://arxiv.org/abs/0707.1835}. See Theorem A.1.
\bibitem{ProveIt}
V. Reshetnikov, \emph{ProveIt},
\path{Analysis/GeneralizedRiemannHypothesis/Article/sections/grp.tex},
commit \texttt{ce35f320bc88c02093e479f5fe5003ff86b2795f}.
\href{https://github.com/VladimirReshetnikov/ProveIt/tree/ce35f320bc88c02093e479f5fe5003ff86b2795f/Analysis/GeneralizedRiemannHypothesis/}{Pinned repository source}.
\end{thebibliography}
\end{document}
